\subsubsection{UC7a – Consulter les points et le niveau (description textuelle complète)}
\begin{longtable}{|p{5cm}|p{11cm}|}
\hline
Cas d'utilisation & Consulter les points et le niveau de fidélité \\ 
\hline
Acteur principal & Utilisateur (Client) \\ 
\hline
Pré-conditions & L'utilisateur est connecté \\ 
\hline
Post-conditions & 
\begin{itemize}
\setlength{\itemsep}{2pt}
\item Le solde de points est affiché
\item Le niveau actuel est affiché
\item La barre de progression vers le niveau suivant est affichée
\item Les coupons disponibles sont listés
\item L'historique des transactions est affiché
\end{itemize} \\ 
\hline
\textbf{Enchaînement nominal} &
\begin{minipage}{\linewidth}
\begin{enumerate}
\setlength{\itemsep}{2pt}
\item L'utilisateur accède à la page "Fidélité".
\item Le système affiche les sections suivantes :
  \begin{itemize}
  \item \textbf{Carte des points} : Points actuels, niveau, barre de progression
  \item \textbf{Collecter des points} : Formulaire pour saisir un code
  \item \textbf{Niveaux} : Présentation des 4 niveaux (Ekoul, Ekoul Kbir, Makhmekh, Makhmekh Kbir)
  \item \textbf{Coupons disponibles} : Récompenses échangeables selon les points
  \item \textbf{Historique} : Transactions passées (acquisition/dépense de points)
  \item \textbf{Comment ça marche} : Explication du programme
  \end{itemize}
\item L'utilisateur peut interagir avec chaque section.
\end{enumerate}
\end{minipage} \\ 
\hline
\textbf{Enchaînements alternatifs} &
\begin{minipage}{\linewidth}
\begin{itemize}
\setlength{\itemsep}{3pt}
\item \textbf{A1 : Saisir un code de fidélité valide}
\begin{enumerate}
\item L'utilisateur saisit un code dans le formulaire "Collecter mes points".
\item L'utilisateur valide le code.
\item Le système vérifie la validité du code.
\item Le système ajoute les points correspondants au compte.
\item Le système met à jour l'affichage des points et du niveau.
\item Le système affiche un message de confirmation.
\item Le système ajoute la transaction à l'historique.
\end{enumerate}
\item \textbf{A2 : Consulter les coupons disponibles}
\begin{enumerate}
\item L'utilisateur consulte la section "Mes coupons disponibles".
\item Le système affiche les coupons éligibles selon le nombre de points :
  \begin{itemize}
  \item 100+ points : Réduction de 10 DT
  \item 250+ points : Entrée offerte
  \item 500+ points : Plat offert
  \item 1000+ points : Menu dégustation
  \end{itemize}
\item Si aucun coupon n'est disponible, le système affiche un message d'encouragement.
\end{enumerate}
\item \textbf{A3 : Consulter l'historique des transactions}
\begin{enumerate}
\item L'utilisateur consulte la section "Historique des points".
\item Le système affiche la liste chronologique des transactions.
\item Chaque transaction montre : description, date/heure, points (+/-).
\item Si l'historique est vide, le système affiche "Aucun historique".
\end{enumerate}
\end{itemize}
\end{minipage} \\ 
\hline
\textbf{Enchaînements d'exception} &
\begin{minipage}{\linewidth}
\begin{itemize}
\setlength{\itemsep}{3pt}
\item \textbf{E1 : Code de fidélité invalide}
\begin{enumerate}
\item L'utilisateur saisit un code invalide.
\item Le système affiche "Code invalide. Veuillez réessayer."
\item Retour à la saisie du code.
\end{enumerate}
\item \textbf{E2 : Code déjà utilisé}
\begin{enumerate}
\item L'utilisateur saisit un code déjà utilisé.
\item Le système affiche "Vous avez déjà utilisé ce code !"
\item Retour à la saisie du code.
\end{enumerate}
\item \textbf{E3 : Champ code vide}
\begin{enumerate}
\item L'utilisateur tente de valider un champ vide.
\item Le système affiche une erreur de validation HTML.
\item Retour à la saisie du code.
\end{enumerate}
\end{itemize}
\end{minipage} \\ 
\hline
\textbf{Notes importantes} &
\begin{minipage}{\linewidth}
\begin{itemize}
\setlength{\itemsep}{2pt}
\item \textbf{Niveaux de fidélité} :
  \begin{itemize}
  \item Ekoul : 0–199 points
  \item Ekoul Kbir : 200–499 points
  \item Makhmekh : 500–999 points
  \item Makhmekh Kbir : 1000+ points
  \end{itemize}
\item \textbf{Barre de progression} : Montre l'avancement vers le prochain niveau.
\item \textbf{Coupons} : Débloqués automatiquement selon le nombre de points.
\item \textbf{Codes de test disponibles} : DINETUNIS100 (100 pts), BIENVENUE50 (50 pts), REPAS2024 (30 pts), FIDELITE20 (20 pts).
\item \textbf{Points totaux gagnés} : Différent des points disponibles (certains peuvent être dépensés).
\item \textbf{Transactions} : Incluent l'acquisition de points via codes et la dépense pour coupons.
\end{itemize}
\end{minipage} \\ 
\hline
\caption{Description textuelle complète – Consulter les points et le niveau de fidélité}
\end{longtable}